\documentclass[11pt,a4paper]{article}
\usepackage[margin=18mm,headheight=16pt,footskip=10mm]{geometry}
\usepackage{fontspec,kotex}
\setmainfont{texgyrepagella-regular.otf}[BoldFont=texgyrepagella-bold.otf,ItalicFont=texgyrepagella-italic.otf,BoldItalicFont=texgyrepagella-bolditalic.otf]
\setmainhangulfont{Apple SD Gothic Neo}[ItalicFont=*,ItalicFeatures={FakeSlant=0.15}]
\setsansfont{texgyreheros-regular.otf}[BoldFont=texgyreheros-bold.otf]
\usepackage{amsmath,amssymb,booktabs,tabularx,array,xcolor,fancyhdr}
\usepackage[hidelinks]{hyperref}
\usepackage{xurl}
\definecolor{ink}{HTML}{173F49}
\pagestyle{fancy}\fancyhf{}
\fancyhead[L]{\small\sffamily\color{ink}THAIE 방법론 비교}
\fancyhead[R]{\small 2026년 11월 19일}
\fancyfoot[C]{\small\thepage\,/\,6}
\setlength{\parindent}{0pt}
\setlength{\parskip}{7pt}
\setlength{\emergencystretch}{3em}
\renewcommand{\arraystretch}{1.35}
\newcolumntype{L}[1]{>{\raggedright\arraybackslash}p{#1}}
\newcommand{\pagetitle}[2]{{\small\sffamily\color{ink}#1\par}\vspace{3mm}{\LARGE\bfseries #2\par}\vspace{4mm}}
\newcommand{\source}[2]{{\small\textbf{출처}\enspace #1\par\url{#2}\par}}
\begin{document}
\pagetitle{01\quad 네 접근법의 측정 대상}{정의, 측정방법과 데이터}
네 가지는 \textbf{측정하려는 대상부터 다릅니다.} Eloundou는 AI로 시간을 줄일 수 있는 과업, Anthropic은 실제 Claude 사용으로 확인되는 업무 범위, METR는 모델이 완수할 수 있는 과업의 길이, THAIE는 과업별 평가 점수에 모델의 시간별 완료 역량을 반영한 값을 측정합니다.

여기서 Anthropic은 Economic Index를 바탕으로 만든 \textbf{2026년 3월의 observed exposure}를 중심으로 비교합니다. 2026년 1월의 \textit{effective AI coverage}도 관련되므로 함께 구분합니다.

{\small
\begin{tabularx}{\textwidth}{L{24mm}XXX}
\toprule
방식 & 정의 & 측정방법 & 데이터\\\midrule
Eloundou et al. (2023) & LLM 활용으로 같은 품질을 유지하며 과업 수행시간을 절반 이상 줄일 수 있는 정도 & 인간과 GPT-4가 과업을 E0·E1·E2 등으로 분류하고, 범주를 수치화해 직업별로 집계 & O*NET 과업 설명, 인간 평가, GPT-4 평가\\\addlinespace
Anthropic observed exposure (2026) & 이론적으로 가능한 과업 중 실제 업무 맥락에서 Claude 활용이 관측되는 범위 & Eloundou 평가, 업무 관련 사용량 기준, 자동화·보조 사용 구분을 결합하고 직업 내 시간 비중으로 집계 & O*NET, Claude.ai 및 API 사용 기록, 이용 방식 분류, 과업별 시간 비중\\\addlinespace
METR & AI 에이전트가 일정 성공률로 완수할 수 있는 과업의 길이. 길이는 인간 전문가 수행시간으로 표시 & 모델이 벤치마크 과업을 실제로 수행한 성공·실패를 인간 수행시간과 연결해 완료 곡선을 추정 & 소프트웨어·머신러닝·보안 과업, 인간 기준시간, 모델별 실행 결과\\\addlinespace
Ours, THAIE & 과업별 AI 활용 평가 점수를 인간 수행시간과 모델별 완료 역량에 따라 가중한 지표 & 기존 평가와 인간 수행시간을 각각 임베딩 기반 릿지 회귀로 한국 과업에 전이하고, METR 완료 곡선을 적용한 뒤 직업별 집계 & Eloundou 평가, Anthropic 시간 추정치, 한국직업사전, 다국어 임베딩, METR 완료 지평\\\bottomrule
\end{tabularx}}

\textbf{표기 안내}\enspace 아래 수식은 비교를 위해 같은 기호 체계로 정리한 것으로, 각 논문의 원래 기호를 그대로 옮긴 것은 아닙니다. 과업(task)은 직무(job)를 구성하는 세부 업무입니다.

\newpage
\pagetitle{02\quad Eloundou et al. (2023)}{평가 범주를 점수로 바꾸는 방식}
LLM을 활용하면 \textbf{동일한 품질을 유지하면서 과업 수행시간을 절반 이상 줄일 수 있는가}를 판단합니다. 실제 과업을 모델에 수행시켜 성공률을 측정하는 방식과는 구별됩니다.

대표적인 $\beta$ 점수는 다음과 같습니다.
\[
b_s=\begin{cases}
1 & \text{LLM을 직접 활용해 수행시간을 절반 이상 단축 가능}\\[4pt]
0.5 & \text{추가 LLM 기반 소프트웨어가 있으면 가능}\\[4pt]
0 & \text{해당 기준에 미달}
\end{cases}
\]
직업별로는 과업 점수를 가중평균합니다.
\[
B_o=\frac{\sum_{s\in S_o}w_{os}b_s}{\sum_{s\in S_o}w_{os}}
\]
$S_o$는 직업 $o$의 과업 집합이고, $w_{os}$는 집계 가중치입니다. 원 연구의 직업별 집계에서는 핵심 과업에 보조 과업보다 높은 가중치를 부여합니다.

원 연구에는 직접 활용만 세는 $\alpha$, 추가 소프트웨어까지 포함하는 $\gamma$ 등의 대안도 있습니다. \textbf{범주값 0.5는 실제 시간 절약률 50\%를 뜻하지 않습니다.}

\textbf{우리 연구와의 연결}\par
우리 코드는 기존 인간의 집계 분류와 GPT-4 분류를 각각 $0$, $0.5$, $1$로 변환해 평균한 값을 학습 대상으로 사용합니다. 이 값을 과업 문장의 임베딩으로 예측하는 릿지 회귀를 학습해 한국 과업에 전이합니다. 한국 과업의 $\hat e_s$는 이 예측값이며, AI 관련 문장과의 코사인 유사도가 아닙니다.

\source{Eloundou et al. (2023), \textit{GPTs are GPTs}, 평가 지침과 측정방법.}{https://arxiv.org/html/2303.10130v5}

\newpage
\pagetitle{03\quad Anthropic observed exposure (2026)}{실제 사용이 확인되는지를 추가하는 방식}
2026년 3월 지표의 구조를 요약하면 다음과 같습니다.
\[
A_o=\sum_{s\in S_o}\tau_{os}
\underbrace{\mathbf1(b_s>0)}_{\text{이론적 가능성}}
\underbrace{\mathbf1(U_s\geq c)}_{\text{충분한 업무 관련 사용}}
\underbrace{a_s}_{\text{자동화 방식 가중치}}
\]
\begin{itemize}
\item $\tau_{os}$는 직업 내 과업의 시간 비중입니다.
\item $U_s$는 Claude.ai 업무 관련 사용과 API 사용을 반영한 사용량입니다.
\item $a_s$는 보조적 활용보다 자동화된 활용에 높은 값을 줍니다.
\end{itemize}
따라서 \textbf{이론적으로 가능해도 사용량 기준을 통과하지 못하면 0}이 됩니다. 반대로 이 지표의 사용 관측은 경제 전체의 도입률을 직접 측정한 것은 아닙니다.

\textbf{2026년 1월 지표와의 구분}\par
별도로 2026년 1월의 \textit{effective AI coverage}는 \textbf{과업별 성공률을 시간·빈도에 따른 비중으로 가중}합니다. 즉, Anthropic도 이미 시간과 성공률을 결합했습니다. 다만 3월 기본 지표에서는 성공률을 제외하고 사용·자동화 중심으로 구성했습니다.

이 구분은 우리 논문의 차별성을 설명할 때 중요합니다. Anthropic의 모든 지표를 단순한 사용 건수 집계로 설명하거나, 시간과 성공 가능성을 고려하지 않았다고 표현하면 안 됩니다.

\source{Anthropic (2026), \textit{Labor market impacts of AI}, 본문과 수리적 부록.}{https://www.anthropic.com/research/labor-market-impacts}
\source{Anthropic (2026), \textit{Economic Index report: Economic primitives}.}{https://www.anthropic.com/research/anthropic-economic-index-january-2026-report}

\newpage
\pagetitle{04\quad METR}{실제 벤치마크 성공으로 완료 곡선을 추정하는 방식}
METR는 모델이 벤치마크 과업을 실제로 수행한 성공·실패를 인간 수행시간과 연결합니다.
\[
P(Y_{sm}=1\mid d_s)=
\left[1+\exp\left\{
\kappa_m\log_2\left(\frac{d_s}{h_m}\right)
\right\}\right]^{-1}
\]
여기서 $Y_{sm}$은 모델 $m$의 과업 $s$에 대한 실제 벤치마크 성공 여부입니다. 이 곡선에서 성공률이 50\%가 되는 인간 수행시간이 $h_m$입니다.

\textbf{$h_m$은 AI가 실제로 작동한 시간이나 모든 과업을 처리할 수 있는 최대시간이 아닙니다.}

\begin{itemize}
\item 과업의 길이는 인간 전문가가 해당 과업을 수행하는 데 필요한 시간으로 표시합니다.
\item 모델별 과업 수행 결과를 이용해 성공률과 인간 수행시간의 관계를 추정합니다.
\item 주요 평가 영역은 소프트웨어 엔지니어링, 머신러닝, 사이버보안입니다.
\item 실제 직장의 맥락이나 모든 직업에 동일한 성공률이 적용된다는 뜻은 아닙니다.
\end{itemize}

\textbf{THAIE와의 관계}\par
METR는 완료 곡선의 근거가 되는 모델 수행 결과를 제공합니다. THAIE는 이 곡선을 한국 과업에 직접 재추정하지 않고, 한국 과업의 예측 인간 수행시간에 적용합니다. 따라서 벤치마크 안에서 추정한 성공확률과 한국 과업에 이전한 완료 계수를 구분해야 합니다.

\source{METR, \textit{Task-Completion Time Horizons of Frontier AI Models}, 공식 방법론.}{https://metr.org/time-horizons/}

\newpage
\pagetitle{05\quad Ours, THAIE}{기존 평가 점수에 이전된 완료 곡선을 적용하는 방식}
먼저 두 값을 별도로 예측합니다.
\[
\hat e_s=\operatorname{clip}_{[0,1]}
\left(\hat{\mathbf w}_e^\top\phi(x_s)+\hat b_e\right),
\qquad
\hat d_s=\exp\left(\hat{\mathbf w}_d^\top\phi(x_s)+\hat b_d\right)
\]
$x_s$는 개별 업무 설명 문장이고, $\phi(x_s)$는 다국어 임베딩입니다. $\hat e_s$는 기존 인간·GPT-4 평가를 전이한 예측 점수이고, $\hat d_s$는 인간 수행시간 예측치입니다. 두 회귀는 서로 다른 학습 대상을 사용합니다.

그다음 모델별 완료 계수를 적용합니다.
\[
\boxed{\lambda_s(m)=\hat e_s\,p_m(\hat d_s)}
\]
\[
p_m(d)=\left[1+\exp\left\{
\kappa_m\log_2\left(\frac{d}{h_m}\right)\right\}\right]^{-1},
\qquad
\kappa_m=\frac{\ln4}{\log_2(h_m/h_{80,m})}
\]
직업별 THAIE는 사전 항목 안에서 과업 점수를 평균하고, 해당 직업의 사전 항목들을 다시 평균합니다.
\[
\Lambda_o(m)=\frac1{|J_o|}\sum_{j\in J_o}
\frac1{|S_j|}\sum_{s\in S_j}\lambda_s(m)
\]
$J_o$는 직업 $o$에 연결된 사전 항목 집합이며, $S_j$는 사전 항목 $j$의 과업 집합입니다.

\textbf{METR와 달리 한국 과업을 모델에 실제로 수행시켜 성공 여부를 측정하지 않았습니다.} 예측 수행시간을 연결고리로 삼아 METR 곡선을 이전한 것입니다.

Anthropic 시간 추정치를 전이해 얻은 $\hat d_s$는 한국 근로자의 실제 평균시간을 관측한 값이 아닙니다. $\lambda_s(m)$ 역시 실제 시간 절약률이나 과업 대체확률이 아닙니다.

\newpage
\pagetitle{06\quad 장점, 한계와 차별성}{무엇을 주장할 수 있는가}
{\small
\begin{tabularx}{\textwidth}{L{25mm}XX}
\toprule
방식 & 장점 & 단점·해석의 한계\\\midrule
Eloundou & 평가 기준이 명확하고 폭넓은 직업을 비교할 수 있음. 실제 도입 이전의 잠재적 활용 범위도 포착 & 판단 기반이며 평가 당시 기술과 가정에 의존. 범주 수치화에 선택이 개입하고 실제 도입·자율 완수·고용효과를 측정하지 않음\\\addlinespace
Anthropic & 실제 사용 여부와 사용 방식을 반영. 이론적 가능성과 활용 사이의 간극을 보여줌 & Claude 이용자·플랫폼의 선택과 도입 여건에 좌우됨. 미관측이 불가능을 뜻하지 않으며, 사용 기록상의 성공도 통제된 벤치마크 성공과 다름\\\addlinespace
METR & 고정된 과업에서 실제 모델 수행을 평가하므로 모델 간 역량 비교에 강함. 성공률과 인간 시간을 연결 & 소프트웨어 중심이고 실험 환경·에이전트 구성에 의존. 직업 전체나 한국 현장 업무로 직접 일반화하기 어려움\\\addlinespace
THAIE & 한국 과업을 동일한 체계로 점수화하고, 모델별 완료 곡선을 바꿔 역량 변화에 따른 차이를 비교할 수 있음 & 기존 평가, 시간 추정, 언어 간 전이, 벤치마크의 직업 간 이전 가정이 누적됨. 실제 도입률·시간 절약률·업무 성공확률은 아님\\\bottomrule
\end{tabularx}}

\textbf{우리 논문의 차별성}\par
THAIE는 기존 과업 평가를 한국 과업 설명에 전이하고, 과업별 인간 수행시간 예측치에 모델별 METR 완료 곡선을 적용하여 직업별 지표를 구성합니다.

따라서 \textbf{‘시간과 성공률을 최초로 결합했다’는 주장은 피해야 합니다.} Anthropic에도 관련 결합이 있습니다. 우리가 강조할 부분은 \textbf{한국 과업으로의 전이, 과업 길이와 모델별 완료 곡선의 명시적 연결, 동일한 과업 체계에서 모델 역량 변화에 따른 비교}입니다.

\textbf{해석 범위}\par
위 장점과 한계는 측정 구조에 근거한 비교입니다. 네 지표를 동일한 데이터와 평가 조건에서 직접 비교해 어느 하나의 예측 성능이 더 우수하다고 입증한 결과는 아닙니다. 이 비교 문서는 기존 분석 수치나 원고의 추정 결과를 변경하지 않습니다.
\end{document}
